% ============================================================================
%  Capitolo 16 -- La memoria permanente                          (Lezione 16)
% ============================================================================
\chapter{La memoria permanente}
\label{ch:storage}

\epigraphline{\ldots un dato non esiste se non è stato salvato\ldots}
  {Prof.\ Almerico Murli, anni Novanta}

\begin{chapgoals}
Alla fine di questo capitolo saprete
\begin{itemize}
\item nominare le unità funzionali di un calcolatore e riconoscerle su una scheda madre e sul kit \prg;
\item distinguere ROM, PROM, EPROM, EEPROM e flash, e dire che cos'è il firmware;
\item spiegare perché una cache rende più veloce un calcolatore, e che cos'è un cache miss;
\item descrivere tracce, settori e blocchi di un disco magnetico e le componenti del suo tempo di accesso;
\item spiegare come un disco ottico memorizza i bit, e confrontare CD e DVD;
\item collocare flash, dischi a stato solido, dischi rigidi e nastri in un sistema di memorizzazione a più livelli.
\end{itemize}
\end{chapgoals}
% Generated by tools/textbook/export_listings.js from tools/decks/mastery.js and tools/concepts/registry.json: do not edit.
\begin{chapstatus}{Stato di apprendimento di questo capitolo}
\masterychip{now}{DA COMPRENDERE OGGI}\enspace ROM, flash, firmware; l'architettura di von Neumann; la cache; la memoria secondaria.\par\smallskip
\masterychip{use}{DA UTILIZZARE OGGI}\enspace il programma completo, gli strumenti di PRG32 (si usa, non si studia, in questo corso).\par\smallskip
\masterychip{later}{DA COMPRENDERE PIÙ AVANTI}\enspace niente.
\end{chapstatus}


% ----------------------------------------------------------------------------
\section{Riscaldamento sulla macchina}

\asmfile[caption={Spegnere e riaccendere il kit.},label=lst:st:warm]{code/l16/c03.S}

Un gioco tiene il punteggio migliore in una variabile
(Listato~\ref{lst:st:warm}). Dopo che il kit è stato spento e riacceso il
titolo è ancora lì, ma il punteggio migliore è 0. La ragione sta nelle due
direttive di sezione. La stringa è in \code{.rodata}, che fa parte del file
della cartuccia, e la cartuccia è un file nella memoria \emph{flash}, che è
non volatile. La variabile è in \code{.data}, che viene copiata in RAM
all'avvio e vive nella SRAM, che è volatile.

Per conservare un punteggio per sempre, bisogna scriverlo in una memoria
permanente. Quali esistono? (Su \prg, i punteggi persistenti sono un servizio
del Cartridge Store; qui il punto è dove vive ciascuna sezione di un
programma.)

% ----------------------------------------------------------------------------
\section{I componenti di un calcolatore}

\subsection{Unità funzionali}

\begin{description}[style=nextline,leftmargin=1.5em]
\item[Processore.] Esegue le istruzioni. Sul kit: il core \riscv
  dell'ESP32-C6.
\item[Memoria centrale.] Contiene i programmi in esecuzione e i loro dati. Sul
  kit: 512~KB di SRAM dentro il chip.
\item[Memoria secondaria.] Conserva programmi e dati quando l'alimentazione è
  spenta. Sul kit: il chip flash che contiene il firmware e le cartucce.
\item[Dispositivi di ingresso/uscita.] Scheda video e monitor, tastiera, mouse,
  stampante, memoria di massa. Sul kit: display, joystick, amplificatore.
\item[Linee di comunicazione.] I bus: linee di dati e di indirizzi condivise da
  tutte le unità.
\end{description}

Un personal computer distribuisce queste unità su una scheda madre; un
microcontrollore le mette quasi tutte in un solo contenitore.

\subsection{La scheda madre}

Una \term{scheda madre} (Figura~\ref{fig:st:board}) porta gli zoccoli per uno
o più processori e, vicino a essi, gli alloggiamenti per i moduli di memoria;
un \emph{chipset} che collega processore, memoria e periferiche; bus di
espansione, oggi alloggiamenti PCI Express per schede video e altre schede;
connessioni per la memoria di massa (Serial ATA, M.2) e per il mondo esterno
(USB, Ethernet, audio); e un chip flash con il firmware che avvia la macchina.
Dal punto di vista architetturale una scheda del 2002 e una di oggi hanno
molto in comune.

\begin{figure}[htb]
\centering
\includegraphics[width=0.7\textwidth]{workstation_board.jpg}
\caption{La scheda madre di un server biprocessore.}
\label{fig:st:board}
\end{figure}

\subsection{Il firmware: il BIOS e i suoi eredi}

\begin{definition}{Firmware}
Il \term{firmware} è il software di base di una macchina, conservato in una
memoria non volatile e usato dal processore dal momento in cui viene acceso.
\end{definition}

In un personal computer si chiamava BIOS: un programma in un chip non volatile,
specifico di una data scheda madre. Le sue istruzioni e i suoi dati sono divisi
in quattro parti: il \emph{POST} (power-on self test), che controlla che i
dispositivi funzionino; il \emph{SETUP}, la configurazione dei dispositivi; il
\emph{BOOT}, che avvia il sistema operativo; e i \emph{DRIVER}, funzioni di
base a basso livello per la tastiera, i dischi e altri dispositivi. I PC di
oggi usano il firmware UEFI, che svolge gli stessi compiti.

\begin{prgmachine}
Sull'ESP32-C6 un programma di avvio in ROM avvia un bootloader in flash, che
avvia il firmware \prg, che carica la cartuccia scelta dal suo menu
\emph{Setup}. Ogni stadio trova, controlla e avvia il successivo: una catena
che va dalla prima istruzione dopo il reset, attraverso strati di memoria
permanente, alla prima istruzione della vostra \code{init}.
\end{prgmachine}

% ----------------------------------------------------------------------------
\section{Memorie non volatili}

\subsection{Dalla ROM alla flash}

Le memorie a semiconduttore che conservano il loro contenuto senza
alimentazione si sono evolute verso un modo sempre più facile di scriverle
(Tabella~\ref{tab:st:rom}). Leggerle richiede un tempo dell'ordine dei
nanosecondi.

\begin{table}[htb]
\centering\small
\caption{Dalla ROM alla flash.}\label{tab:st:rom}
\begin{tabularx}{\textwidth}{@{}lP{0.24\textwidth}LL@{}}
\toprule
\textbf{Memoria} & \textbf{Scritta da} & \textbf{Cancellata} & \textbf{Uso tipico}\\
\midrule
ROM & il costruttore, una volta per tutte & mai & codice di avvio, vecchie cartucce di gioco\\
PROM & l'utente, una volta & mai & piccole serie di produzione\\
EPROM & l'utente, molte volte & tutta insieme, con luce ultravioletta attraverso una finestra di quarzo (circa 15 minuti) & sviluppo, BIOS dei primi PC\\
EEPROM & l'utente, sul posto & elettricamente, un byte alla volta & dati di configurazione\\
Flash & l'utente, sul posto & elettricamente, a blocchi & firmware, chiavette USB, schede di memoria, dischi a stato solido\\
\bottomrule
\end{tabularx}
\end{table}

Le sigle stanno per \emph{read-only memory} (memoria di sola lettura), ROM
\emph{programmabile}, PROM \emph{cancellabile} e PROM \emph{cancellabile
elettricamente}. Una ROM è, dal punto di vista logico, una macchina
combinatoria: una tabella che associa una parola fissa a ogni indirizzo, come
fanno le tabelle \code{dx} e \code{dy} di \vermone.

\subsection{La memoria flash}

La flash è il tipo più recente di EEPROM: può essere riprogrammata sul posto,
senza essere rimossa dalla scheda, ed è oggi la memoria a semiconduttore meno
costosa. Ha tre peculiarità.

\begin{itemize}
\item Si cancella a \emph{blocchi} di byte, non un byte alla volta. La
  cancellazione pone a 1 tutti i bit di un blocco, e la scrittura può solo
  trasformare degli 1 in 0.
\item Ogni blocco sopravvive a un numero limitato di cicli di cancellazione,
  così che il controller distribuisce le scritture in modo uniforme sui
  blocchi.
\item La lettura è veloce; la scrittura e, soprattutto, la cancellazione sono
  lente.
\end{itemize}

La flash del kit contiene il firmware e le cartucce: caricare una cartuccia
cancella e riscrive alcuni dei suoi blocchi. Ed è per questo che il punteggio
migliore del riscaldamento non viene semplicemente riscritto nella cartuccia a
ogni frame: la flash non è un posto per le variabili.

\begin{figure}[htb]
\centering
\begin{subfigure}[b]{0.3\textwidth}\centering\includegraphics[width=\textwidth,height=34mm,keepaspectratio]{eprom.jpg}\caption{Una EPROM con la sua finestra di quarzo}\end{subfigure}\hfill
\begin{subfigure}[b]{0.3\textwidth}\centering\includegraphics[width=\textwidth,height=34mm,keepaspectratio]{hard_disk.jpg}\caption{Un disco rigido aperto}\label{fig:st:hdd}\end{subfigure}\hfill
\begin{subfigure}[b]{0.3\textwidth}\centering\includegraphics[width=\textwidth,height=34mm,keepaspectratio]{ssd.jpg}\caption{Un disco a stato solido}\end{subfigure}
\caption{Tre memorie permanenti.}
\end{figure}

% ----------------------------------------------------------------------------
\section{La cache}

\subsection{Perché una cache?}

Durante i cicli di clock necessari per leggere o scrivere la memoria il
processore aspetta; più lenta è la memoria, più cicli si perdono. Perché non
costruire memorie veloci quanto i registri del processore? Perché dovrebbero
stare sullo stesso chip, dato che un collegamento attraverso un bus è più
lento; e allora non potrebbero essere molto grandi, o il chip crescerebbe in
dimensioni e in costo. Il compromesso è una memoria piccola e veloce sul chip
del processore, e una o più memorie più grandi e più lente raggiunte
attraverso il bus.

\subsection{Come funziona una cache}

\begin{definition}{Cache}
Una \term{cache} (dal francese \emph{cacher}, nascondere) è una memoria di
piccola capacità ma veloce, posta tra il processore e la memoria centrale, che
tiene una copia di parte del contenuto di quest'ultima.
\end{definition}

Il processore cerca prima nella cache (Figura~\ref{fig:st:cache}). C'è un
\term{successo} \eng{hit} se il contenuto dell'indirizzo richiesto è lì, e
arriva subito. C'è un \term{fallimento} \eng{miss} se non c'è, e allora deve
essere copiato dalla RAM nella cache; con esso viene copiato un numero fisso
di locazioni vicine. Il programma non si ferma e non se ne accorge: l'accesso
è semplicemente più lento.

\begin{figure}[htb]
\centering
\begin{tikzpicture}
  \node[blkacc,minimum width=36mm,minimum height=9mm] (p) at (0,0) {processore};
  \node[blkamb,minimum width=36mm,minimum height=11mm] (c) at (0,-1.8) {\textbf{cache}\\piccola, veloce};
  \node[blkvio,minimum width=36mm,minimum height=13mm] (m) at (0,-3.9) {\textbf{memoria centrale}\\grande, lenta};
  \draw[<->,bus] (p) -- (c) node[midway,right,lbl]{ogni accesso};
  \draw[<->,bus] (c) -- (m) node[midway,right,lbl]{solo su un miss: un blocco};
\end{tikzpicture}
\caption{Una cache tra il processore e la memoria centrale.}
\label{fig:st:cache}
\end{figure}

Perché copiare i vicini? Per una proprietà empirica dei programmi.

\begin{prgconcept}[Il principio di località]
Un programma che ha appena usato un indirizzo probabilmente lo userà di nuovo
presto (\emph{località temporale}) e probabilmente userà i suoi vicini
(\emph{località spaziale}). I cicli e gli array ne sono la ragione.
\end{prgconcept}

Più successi ci sono, più breve è il tempo medio di accesso. Se una frazione
$h$ degli accessi sono successi, la cache risponde in un tempo $t_c$ e un
fallimento costa in più il tempo $t_m$ della memoria, la media è
\[ t = h\,t_c+(1-h)\,(t_c+t_m). \]
Con $t_c=1$~ns, $t_m=50$~ns e il 95\% di successi,
$t=0,95\times 1+0,05\times 51=3,5$~ns: il sistema si comporta quasi come se
l'intera memoria fosse veloce quanto la cache.

\subsection{Livelli di cache e macchine tipiche}

I processori reali hanno più livelli di cache, tutti costruiti con SRAM
(Tabella~\ref{tab:st:levels}). Il primo livello, interno, mantiene un flusso
continuo dentro il processore; il secondo limita il rallentamento dovuto alla
memoria centrale; il terzo, da 8 a 32~MB e oltre, è condiviso dai core dello
stesso processore.

\begin{table}[htb]
\centering\small
\caption{I livelli di cache su tre macchine.}\label{tab:st:levels}
\begin{tabularx}{\textwidth}{@{}lllL@{}}
\toprule
 & \textbf{PC del 2002} & \textbf{PC del 2020} & \textbf{ESP32-C6 (il kit)}\\
\midrule
Processore & 1 core & 4 core & 1 core \riscv, 160 MHz\\
Cache L1 & 64 KB & 64 KB per core & nessuna per la SRAM\\
Cache L2 & 512 KB & 512 KB per core & nessuna\\
Cache L3 & nessuna & 8 MB, condivisa dai core & nessuna\\
Memoria centrale & 256 MB & 16 GB & 512 KB di SRAM\\
Memoria di massa & disco rigido da 40 GB & 1 TB + 1 TB, in mirror & pochi MB di flash, letta attraverso una piccola cache\\
\bottomrule
\end{tabularx}
\end{table}

Sull'ESP32-C6 la SRAM è già veloce quanto serve al core; solo la flash esterna
sta dietro una cache.

% ----------------------------------------------------------------------------
\section{La memoria secondaria}

\subsection{Caratteristiche}

La \term{memoria secondaria}, o \term{memoria di massa}, conserva tutti i
programmi e i dati di un calcolatore, mentre la memoria centrale contiene i
programmi in esecuzione e i dati di cui hanno bisogno. Rispetto alla memoria
centrale è

\begin{itemize}
\item non volatile: i dati non si perdono quando il calcolatore viene spento;
\item di grande capacità: più ordini di grandezza in più;
\item di basso costo per bit: più ordini di grandezza in meno;
\item più lenta: di alcuni ordini di grandezza.
\end{itemize}

La sua informazione è organizzata in \term{blocchi}, di solito di un kilobyte
o più, non in singoli byte. Alcuni supporti permettono l'accesso
\emph{diretto} a qualsiasi blocco; altri solo l'accesso \emph{sequenziale}.

\subsection{Dischi magnetici: tracce, settori, blocchi}

Un disco magnetico (Figura~\ref{fig:st:hdd}) è fatto di piatti rivestiti di un
materiale magnetizzabile e tenuti in rotazione a velocità costante, e di
testine che leggono e scrivono l'informazione. La superficie è divisa in
\term{tracce} concentriche e in \term{settori}, spicchi del disco
(Figura~\ref{fig:st:disk}). Un settore taglia ogni traccia in blocchi: il
blocco è la più piccola unità indirizzabile, identificata dal numero della sua
traccia e del suo settore. \term{Formattare} un disco significa dividerne la
superficie in tracce e settori. L'accesso è diretto: la testina può essere
posizionata su qualsiasi blocco senza passare su tutti i precedenti. Più
settori formano un \emph{cluster}; le tracce allo stesso raggio su tutti i
piatti formano un \emph{cilindro}.

\begin{figure}[htb]
\centering
\begin{tikzpicture}
  \fill[prgpanel] (0,0) circle (2.6);
  \fill[prgamber!45] (0,0) -- (30:2.6) arc (30:60:2.6) -- cycle;
  \begin{scope}
    \clip (0,0) -- (30:2.7) arc (30:60:2.7) -- cycle;
    \fill[prgred!70,even odd rule] (0,0) circle (1.9) (0,0) circle (1.5);
  \end{scope}
  \draw[prgaccent,line width=2pt] (0,0) circle (1.7);
  \foreach \r in {0.7,1.1,1.5,1.9,2.3,2.6} \draw[prgsteel,thin] (0,0) circle (\r);
  \foreach \a in {0,30,...,330} \draw[prgsteel,thin] (\a:0.7) -- (\a:2.6);
  \fill[white] (0,0) circle (0.7); \draw[prgsteel] (0,0) circle (0.7); \fill[prgsteel] (0,0) circle (0.12);
  \draw[<-,thick] (150:1.7) -- (-3.6,1.6) node[left,lbl,text=prgaccent]{una traccia};
  \draw[<-,thick] (45:2.4) -- (3.4,2.2) node[right,lbl,text=prgamber]{un settore};
  \draw[<-,thick] (45:1.7) -- (3.4,0.6) node[right,lbl,text=prgred]{un blocco};
  \draw[line width=3pt,prgink] (3.3,-2.4) -- (-10:1.75);
  \fill[prgink] (-10:1.75) circle (0.1);
  \node[lbl,anchor=west] at (3.4,-2.4) {braccio e testina};
\end{tikzpicture}
\caption{La superficie di un disco magnetico.}
\label{fig:st:disk}
\end{figure}

\subsection{Tempo di accesso di un disco magnetico}

Per leggere o scrivere un blocco la testina deve essere portata all'indirizzo
giusto; durante il trasferimento la testina è ferma e il disco si muove. Il
tempo di accesso è la somma di tre componenti.

\begin{description}[style=nextline,leftmargin=1.5em]
\item[Tempo di posizionamento.] Il tempo per spostare la testina in senso
  radiale fino al cilindro giusto, e quindi alla traccia giusta.
\item[Tempo di latenza.] Il tempo di attesa perché il settore desiderato passi
  sotto la testina; in media, mezzo giro.
\item[Tempo di trasferimento.] Il tempo perché il blocco passi sotto la
  testina.
\end{description}

Un disco che gira a 7\,200 giri al minuto impiega $60/7200 = 8,33$~ms per
giro, e la sua latenza media è di circa 4,2~ms. I dischi rigidi sono sigillati
contro la polvere, hanno capacità di terabyte e tempi di accesso di
millisecondi. Un millisecondo è un milione di nanosecondi: durante un solo
accesso a un disco un processore a 160~MHz potrebbe eseguire più di centomila
istruzioni.

\subsection{Dischi rigidi e dischi a stato solido}

Un \term{disco rigido} (HDD, \emph{hard disk drive}) ha piatti magnetici e
testine mobili; il suo tempo di accesso, di millisecondi, è dominato dalla
meccanica; ha il costo più basso per terabyte ed è ancora la scelta per i
grandi archivi e per i backup. Un \term{disco a stato solido} (SSD) è memoria
flash con un controller, e non ha parti mobili; il suo tempo di accesso è di
decine di microsecondi; si collega attraverso SATA o, più velocemente,
attraverso NVMe su PCI Express. È la memoria di massa standard di portatili e
telefoni, ed è la stessa tecnologia del chip flash del kit.

% ----------------------------------------------------------------------------
\section{Dischi ottici}

\subsection{Bit letti con la luce}

Un disco ottico è un disco rigido sul quale viene proiettato un raggio laser
di pochi milliwatt; la riflessione cambia tra i \term{pit} (le fossette) e la
superficie, il \term{land} (Figura~\ref{fig:st:pits}).

\begin{prgwarn}
0 e 1 \emph{non} sono pit e land. Un 1 è un \emph{cambiamento}, da pit a land o
da land a pit; uno 0 è nessun cambiamento.
\end{prgwarn}

Codificare un 1 come una transizione è un'idea usata anche in molti
collegamenti seriali, perché permette al ricevitore di recuperare il ritmo dei
bit dal segnale stesso. I dischi ottici di solito si scrivono una volta sola,
perché i cambiamenti fisici sono irreversibili: sono usati per distribuire
programmi e come archivi che non devono essere modificati. In un
CD-Recordable, due strati di policarbonato racchiudono un sottile foglio
riflettente coperto da un colorante traslucido, e la registrazione brucia dei
punti nel colorante; i dischi riscrivibili usano un materiale che può tornare
indietro.

\begin{figure}[htb]
\centering
\begin{subfigure}[b]{0.3\textwidth}\centering\includegraphics[width=\textwidth,height=34mm,keepaspectratio]{pit_land.jpg}\caption{Pit e land (microscopio elettronico)}\label{fig:st:pits}\end{subfigure}\hfill
\begin{subfigure}[b]{0.3\textwidth}\centering\includegraphics[width=\textwidth,height=34mm,keepaspectratio]{cdrom_drive.jpg}\caption{Dentro un lettore di CD-ROM}\end{subfigure}\hfill
\begin{subfigure}[b]{0.3\textwidth}\centering\includegraphics[width=\textwidth,height=34mm,keepaspectratio]{dvd_drive.jpg}\caption{Dentro un lettore di DVD}\end{subfigure}
\caption{Memorie ottiche.}
\end{figure}

\subsection{Il compact disc}

Su un compact disc l'informazione è memorizzata su un'unica traccia a spirale,
con densità di registrazione costante; la distanza tra due giri successivi
della spirale, il \emph{passo della traccia}, è di 1,6 micrometri. Sono
possibili due modi di far girare il disco. Con la \emph{velocità lineare
costante} (CLV) il disco gira più in fretta quando la testina legge vicino al
centro, da 200 a 500 giri al minuto, così che la traccia passa sotto la testina
sempre alla stessa velocità (1,2~m/s alla velocità di base, fissata dalla
velocità di bit necessaria per riprodurre la musica). Con la \emph{velocità
angolare costante} (CAV) i giri sono costanti e così il numero di bit per
settore, con il risultato che la superficie esterna è sottoutilizzata.

\subsection{DVD e CD a confronto}

Il DVD (\emph{digital versatile disk}) ha pit più piccoli e una spirale più
fitta, e fino a due strati e due facce (Tabella~\ref{tab:st:dvd}).

\begin{table}[htb]
\centering\small
\caption{DVD e CD.}\label{tab:st:dvd}
\begin{tabular}{@{}lll@{}}
\toprule
 & \textbf{DVD} & \textbf{CD}\\
\midrule
Diametro del disco & 120 mm & 120 mm\\
Spessore del disco & 1,2 mm ($0,6$ mm $\times$ 2) & 1,2 mm\\
Facce / strati & 1 o 2 / 1 o 2 & 1 / 1\\
Dimensione minima dei pit & 0,4 $\mu$m & 0,83 $\mu$m\\
Passo della traccia & 0,74 $\mu$m & 1,6 $\mu$m\\
Velocità media di trasferimento & 4,7 MB/s & 0,15 MB/s\\
Capacità, 1 strato, 1 faccia & circa 5 GB & circa 0,68 GB\\
Capacità, 2 strati, 2 facce & 17 GB & --\\
\bottomrule
\end{tabular}
\end{table}

% ----------------------------------------------------------------------------
\section{La memorizzazione su più livelli}

Un moderno sistema di calcolo ha più livelli di memoria permanente
(Figura~\ref{fig:st:levels}), spesso gestiti automaticamente: i dati usati
spesso migrano verso i livelli veloci, e il resto scende verso quelli
economici. È l'idea della cache, applicata alla memoria di massa; e da un chip
di pochi megabyte sul kit a sale di petabyte (Figura~\ref{fig:st:racks}) la
gerarchia è la stessa a ogni scala.

\begin{figure}[htb]
\centering
\begin{tikzpicture}[node distance=4mm]
  \node[blkred,text width=21mm,minimum height=17mm] (a) {\textbf{Disco RAM}\\[1pt]\footnotesize il più veloce, volatile};
  \node[blkamb,text width=21mm,minimum height=17mm,right=of a] (b) {\textbf{Disco a stato solido}\\[1pt]\footnotesize di lavoro};
  \node[blkacc,text width=21mm,minimum height=17mm,right=of b] (c) {\textbf{Disco rigido}\\[1pt]\footnotesize capiente, economico};
  \node[blk,text width=21mm,minimum height=17mm,right=of c] (d) {\textbf{Libreria di dischi}\\[1pt]\footnotesize accesi al bisogno};
  \node[blkvio,text width=21mm,minimum height=17mm,right=of d] (e) {\textbf{Libreria a nastri}\\[1pt]\footnotesize archivi};
  \foreach \x/\y in {a/b,b/c,c/d,d/e} \draw[<->,thick] (\x) -- (\y);
\end{tikzpicture}
\caption{Una memoria permanente a più livelli.}
\label{fig:st:levels}
\end{figure}

\begin{figure}[htb]
\centering
\includegraphics[width=0.6\textwidth]{storage_racks.jpg}
\caption{Una libreria di nastri in un centro di calcolo.}
\label{fig:st:racks}
\end{figure}

% ----------------------------------------------------------------------------
% Generated by tools/textbook/export_listings.js from decks/plan.js and code/l16.S: do not edit.
\section{Il programma di questo capitolo}\label{sec:lab:l16}

I listati di questo capitolo sono parti di un solo programma,
\fname{code/l16.S} nel repository del corso, stampato per intero nel
Listato~\ref{lst:lab:l16}. I test automatici del materiale del corso lo
compilano e lo eseguono sull'esecutore senza interfaccia di \prgqt{}, con
ingressi programmati, e controllano il frame; la
Figura~\ref{fig:lab:l16} è il frame scritto dal test.

\begin{description}[style=nextline,leftmargin=1.5em]
\item[Sullo schermo.] Il titolo HIGH SCORE e due barre: il punteggio di questa prova, che cresce mentre A è premuto e torna a zero quando A è rilasciato, e il punteggio migliore, che non scende mai. Riavviate la cartuccia: il titolo c'è ancora, il punteggio migliore è di nuovo zero.
\item[Che cosa premere.] Tenete premuto A, rilasciatelo, poi riavviate la cartuccia.
\end{description}

\begin{figure}[htb]
\centering
\includegraphics[width=0.5\textwidth]{lab_l16.png}
\caption{\fname{l16.S} su \prgqt{}.}\label{fig:lab:l16}
\end{figure}

\asmfile[caption={Il programma completo del capitolo, \fname{code/l16.S}.},label=lst:lab:l16]{code/l16/program.S}

Per eseguirlo su \prgqt{} i comandi sono quelli del
Listato~\ref{lst:lab:l16:run}; il Paragrafo~\ref{sec:prg:testing} dà gli
stessi passi per QEMU e per la scheda.

\begin{shell}[caption={Compilare \fname{l16.S} ed eseguirlo su \prgqt{}.},label=lst:lab:l16:run]
# 1. start PRG32-QT on this computer (it listens on port 8080)

# 2. in the PRG32 folder, with ESP-IDF active: build the cartridge
python3 -m prg32 cartridge build ~/arc12/code/l16.S --portable \
    --entry-prefix l16 --name l16 --out l16.prg32

# 3. send it to PRG32-QT and start it
python3 -m prg32 esp32c6 upload-and-run l16.prg32 \
    --url http://localhost:8080 --slot cart0

# 4. look, press, compare with the previous slide
\end{shell}


% ----------------------------------------------------------------------------
% Generated by tools/textbook/export_listings.js from tools/concepts/assessments.json: do not edit.
\begin{chapstatus}{Il filo del capitolo}
La memoria è una gerarchia: piccola e veloce vicino al processore, grande e lenta lontano, e la località fa funzionare il trucco. Prossimo passo: la terza parte di Verm1, che ha bisogno di memoria per un corpo intero.
\par\medskip\textbf{Biglietto d'uscita} (da due a quattro minuti, da soli, su carta):
\begin{enumerate}[nosep,leftmargin=2em]
\item Che cosa succede in caso di fallimento di cache? Perché un ciclo che legge un array in ordine ha pochi fallimenti?
\item Perché il firmware e le cartucce del kit non possono stare solo in RAM? Quale memoria li conserva, e qual è il suo limite?
\end{enumerate}
\end{chapstatus}


\section{Verificate di aver capito}

\begin{quiz}{Quale memoria si può cancellare elettricamente, a blocchi, senza rimuoverla dalla scheda?}
\item ROM
\item PROM
\item EPROM
\item Flash
\end{quiz}

\begin{quiz}{Che cosa accade su un cache miss?}
\item Il programma si ferma con un errore
\item Il dato viene prelevato dalla memoria centrale e copiato nella cache, con i suoi vicini
\item La cache viene cancellata
\item Il processore salta l'istruzione
\end{quiz}

\begin{quiz}{Quale componente del tempo di accesso di un disco rigido dipende dalla velocità di rotazione?}
\item Il tempo di posizionamento
\item La latenza
\item Il tempo di formattazione
\item Il tempo di avvio
\end{quiz}

\begin{quiz}{Su un disco ottico, un 1 è codificato come:}
\item un pit
\item un land
\item una transizione tra pit e land
\item un doppio pit
\end{quiz}

\section{Esercizi}

\begin{exercise}[su carta e su \prg]
\begin{enumerate}[nosep]
\item Un disco gira a 5\,400 giri al minuto: qual è la sua latenza media?
\item Una cache risponde in 2~ns e la memoria in 80~ns: qual è il tempo medio
  di accesso con il 90\% di successi? E con il 99\%?
\item Dite dove vive sul kit \prg ciascuna di queste cose: il codice di avvio,
  il firmware, una cartuccia, la variabile \code{counter}, il registro
  \reg{t0}.
\item Trovate le specifiche della memoria di massa del vostro portatile o del
  vostro telefono: tecnologia, capacità, interfaccia.
\item Su \prg, cronometrate due cicli che sommano un array grande: uno in
  ordine, uno che salta con un passo grande. C'è una differenza, e perché su
  questo chip potrebbe non esserci?
\end{enumerate}
\expected{Latenza: metà di $60/5400=11,1$~ms, cioè circa 5,6~ms. Tempo di
accesso: $0,9\times 2+0,1\times 82=10$~ns, e
$0,99\times 2+0,01\times 82=2,8$~ns. Le cinque cose vivono in ROM, in flash,
in flash (o in RAM una volta caricata), in SRAM, e in un registro di flip-flop
dentro il processore. Nell'ultimo compito la SRAM dell'ESP32-C6 non ha una
cache davanti, così che lo schema di accesso conta poco; su un portatile la
differenza è grande.}
\end{exercise}

\section{Glossario}

\begin{glossarytable}
motherboard & scheda madre\\
firmware & firmware\\
non-volatile & non volatile\\
cache hit / miss & successo / fallimento di cache\\
mass storage & memoria di massa\\
track / sector / block & traccia / settore / blocco\\
seek time & tempo di posizionamento\\
latency & tempo di latenza\\
solid-state disk & disco a stato solido\\
tape library & libreria a nastri\\
\end{glossarytable}

\section{Conclusioni e letture}

Le memorie non volatili contengono di solito il software di base di una
macchina, il firmware. La cache migliora le prestazioni tenendo vicino al
processore una copia di parte della memoria centrale, e la stessa idea, un
livello piccolo e veloce davanti a uno grande e lento, è usata di nuovo nei
sistemi di memorizzazione. Le memorie permanenti possono essere a
semiconduttore (flash), magnetiche o ottiche, e un sistema moderno ne ha più
livelli, spesso gestiti automaticamente.

\paragraph{Letture.} \citet{patterson2020}, paragrafi 5.1--5.4;
\citet{stallings2021}, i capitoli sulla cache e sulla memoria esterna;
\citet{tanenbaum2013}, capitolo~2.

Il terzo blocco è completo. Il prossimo capitolo lo usa nella terza parte di
\vermone: un corpo, regole e stati.
